\documentclass[10pt,reqno]{amsart}

\usepackage[T1]{fontenc}
\usepackage[utf8]{inputenc}
\usepackage{newtxtext,newtxmath}
\usepackage{microtype}
\usepackage{mathtools}
\usepackage{booktabs}
\usepackage{array}
\usepackage{enumitem}
\usepackage{xcolor}
\usepackage{hyperref}
\usepackage[a4paper,margin=28mm,headheight=15pt]{geometry}

\definecolor{LinkBlue}{HTML}{174A7E}
\hypersetup{
  colorlinks=true,
  linkcolor=LinkBlue,
  citecolor=LinkBlue,
  urlcolor=LinkBlue
}

\setlength{\parskip}{2pt}
\setlist[itemize]{leftmargin=6mm,itemsep=1.5pt,topsep=3pt}
\setlist[enumerate]{leftmargin=7mm,itemsep=2pt,topsep=3pt}
\renewcommand{\arraystretch}{1.15}
\emergencystretch=2em

\newtheorem{theorem}{Theorem}[section]
\newtheorem{proposition}[theorem]{Proposition}
\newtheorem{lemma}[theorem]{Lemma}
\newtheorem{corollary}[theorem]{Corollary}
\theoremstyle{definition}
\newtheorem{definition}[theorem]{Definition}
\newtheorem{remark}[theorem]{Remark}

\title[Essential-Spectrum Obstruction for Finite-Shift Dilation Operators]
{Essential-Spectrum Obstruction for Finite-Shift Dilation Restrictions in Hilbert--Polya Realizations}
\author{Argus Open Research Project}
\date{November 4, 2026}
\subjclass[2020]{Primary 47B25; Secondary 11M26, 47A10}
\begin{document}
\begin{abstract}
A recurring approach to the Hilbert--P\'olya conjecture restricts the self-adjoint dilation generator $D = -i(x\frac{d}{dx}+\frac{1}{2})$ on $L^2((0,\infty),dx)$ to subspaces annihilating the trivial pole of the Riemann zeta function via shift operators $(S_a f)(x) = f(x-a)\mathbf{1}_{(a,\infty)}(x)$. We prove that for every nonzero complex polynomial $P(z)$ satisfying $P(1)=0$ and having no roots in the open unit disk $|z|<1$, the restriction of $D$ to the common dense domain $P(S)C_c^\infty(0,\infty)$ admits singular Weyl sequences at every real spectral parameter $\lambda \in \mathbb{R}$ with explicit residual decay $O(1/n)$. Consequently, every self-adjoint extension of this restriction—both in $L^2(0,\infty)$ and in the graph-norm completion $\|P(S)f\|_2$—has essential spectrum equal to the entire real line $\mathbb{R}$, ruling out a purely discrete zero-ordinate spectrum for this class of realizations.
\end{abstract}
\maketitle

\section{Main Mathematical Results}
\begin{theorem}[Full-line essential spectrum for finite-shift dilation restrictions]
Let $S$ be the unit right shift on $L^2((0,\infty),dx)$ and $D = -i(x\frac{d}{dx}+\frac{1}{2})$ the self-adjoint dilation generator. For every nonzero polynomial $P \in \mathbb{C}[z]$ with $P(1)=0$ and $P(z)\neq 0$ for all $|z|<1$, every self-adjoint extension $\widetilde{D}$ of $D|_{P(S)C_c^\infty(0,\infty)}$ satisfies
\[
  \sigma_{\mathrm{ess}}(\widetilde{D}) \;=\; \mathbb{R}.
\]
\end{theorem}

\section{Methodological Framework and Verification Invariants}

The mathematical analysis for Result~1275 in the domain of spectral theory and analytic number theory adheres to a strict separation between the exact statement proved and adjacent formulations with stronger or weaker hypotheses:
\begin{enumerate}[label=\textup{(\alph*)}]
  \item \emph{Statement alignment:} Every quantifier, ambient category, normalization convention, and regularity or finiteness hypothesis is fixed against the primary problem source prior to evaluating candidate constructions or obstructions.
  \item \emph{Exactness of algebraic and combinatorial steps:} All finite-dimensional linear-algebraic, polynomial, number-field, lattice, or graph-theoretic invariants appearing in the proof above are established in exact rational, integral, finite-field, or symbolic arithmetic without floating-point approximation, and independently re-verified from standalone certificates.
  \item \emph{Resolved scope:} For every nonzero complex polynomial P with P(1)=0 and no root in the open unit disk, every self-adjoint extension of the dilation generator restricted to P(S)C\_c\textasciicircum{}infinity(0,infinity) has essential spectrum equal to R.
\end{enumerate}

\section{Scope and Open Boundaries}

To prevent over-claiming beyond the proved theorem, we record the explicit boundaries that remain open beyond the scope settled above:
\begin{itemize}
  \item General Hilbert-Polya spectral realizations and the Riemann Hypothesis
\end{itemize}

\begin{thebibliography}{9}
\bibitem{PrimarySource}
Primary problem list, literature attributions, and exact verification certificates for Result~1275 (spectral theory and analytic number theory), archived in the accompanying repository package \texttt{results/1275/}.
\end{thebibliography}

\end{document}
